\documentclass{article}
\usepackage[utf8]{inputenc}
\usepackage{amsmath}
\usepackage{amsfonts}
\usepackage{amssymb}
\usepackage{geometry}
\usepackage{tcolorbox}
\usepackage{listings}
\usepackage{xcolor}
\usepackage{hyperref}

\geometry{a4paper, margin=1in}

\definecolor{codegreen}{rgb}{0,0.6,0}
\definecolor{codegray}{rgb}{0.5,0.5,0.5}
\definecolor{codepurple}{rgb}{0.58,0,0.82}
\definecolor{backcolour}{rgb}{0.95,0.95,0.92}

\lstset{
    backgroundcolor=\color{backcolour},   
    commentstyle=\color{codegreen},
    keywordstyle=\color{magenta},
    numberstyle=\tiny\color{codegray},
    stringstyle=\color{codepurple},
    basicstyle=\ttfamily\small,
    breakatwhitespace=false,         
    breaklines=true,                 
    captionpos=b,                    
    keepspaces=true,                 
    numbers=left,                    
    numbersep=5pt,                  
    showspaces=false,                
    showstringspaces=false,
    showtabs=false,                  
    tabsize=2
}

\title{Module 04: Feature Engineering II - Group-Level Transforms and Contextual Modeling}
\author{Cybernauts 2026 Datathon Campaign}
\date{May 2026}

\begin{document}

\maketitle

\begin{abstract}
Numbers in isolation are often meaningless. If a column says a user spent \$150, is that high or low? Without context, the model cannot know. This note covers the use of the \texttt{.transform()} function to build contextual features that compare individual data points against their peer groups, transforming raw numbers into behavioral signals.
\end{abstract}

\tableofcontents
\newpage

\section{The Concept of Contextual Modeling}
Machine learning models struggle to understand "normality." To help them, we must explicitly calculate what is normal for a specific group and then show the model how an individual row deviates from that norm.

\section{The \texttt{.transform()} Function}
Unlike \texttt{.agg()}, which reduces the number of rows to a summary table, \texttt{.transform()} calculates a statistic for a group but \textbf{broadcasts it back} to the original shape of the DataFrame.

\begin{tcolorbox}[colback=blue!5!white,colframe=blue!75!black,title=Essential Syntax]
\texttt{df['group\_mean'] = df.groupby('Group\_Col')['Target\_Col'].transform('mean')}
\end{tcolorbox}

\section{Practical Contextual Features}

\subsection{Relative Deviation}
Is this user spending more than their personal average? 
\begin{equation}
    \text{Diff\_From\_Avg} = \text{Current\_Spend} - \text{User\_Avg\_Spend}
\end{equation}

\begin{lstlisting}[language=Python]
df['user_avg'] = df.groupby('User_ID')['Spend'].transform('mean')
df['spend_vs_user_avg'] = df['Spend'] - df['user_avg']
\end{lstlisting}

\subsection{Group-Based Z-Scores}
Is this house priced abnormally high for its neighborhood?
\begin{equation}
    z_{group} = \frac{x - \mu_{group}}{\sigma_{group}}
\end{equation}

\begin{lstlisting}[language=Python]
df['hood_avg'] = df.groupby('Neighborhood')['Price'].transform('mean')
df['hood_std'] = df.groupby('Neighborhood')['Price'].transform('std')
df['price_z_hood'] = (df['Price'] - df['hood_avg']) / df['hood_std']
\end{lstlisting}

\section{Practice Problems}
\textbf{P1:} When should you use \texttt{.transform()} instead of \texttt{.merge()}? \\
\textit{Answer: Use \texttt{.transform()} for simple statistical broadcasts (mean, sum, std) because it is faster and requires less code. Use \texttt{.merge()} when you need to join complex summary tables with multiple different metrics.}

\textbf{P2:} You are predicting fraud in transactions. Why is \texttt{Amount - Avg\_Amount\_For\_That\_Card} a better feature than just \texttt{Amount}? \\
\textit{Answer: Because \$500 might be normal for a wealthy user but highly suspicious for a student. The relative difference provides the "behavioral context" that raw numbers lack.}

\end{document}
